% Compounding period exercises
% Stu Alden
% ------------------------------------------
\documentclass[addpoints,twoside]{exam}         % Can't specify font size here, 
                                                % so we'll set it manually below
\usepackage{problemset}                         % Custom package for problem sets - 
                                                % see problemset.sty    
% -------------------------------------------------------------------------
\begin{document}

\FontEasyRead                                   % Set font size here

\Heading                                        % Modify heading as needed
{Financial Literacy}                            % Course
{Compound Interest II - Exercise Set 1}          % Title
{Answer the questions below.}

\ProblemSection{Fill in the chart:}

\renewcommand{\arraystretch}{2} % Increases row height, looks better
\setlength{\tabcolsep}{10pt} % Adjusts column spacing
\begin{table}[h]
    \centering
    \begin{tabular}{|p{3cm}<{\raggedright}|p{3cm}<{\raggedright}|p{2cm}<{\raggedright}|p{5cm}<{\raggedright}|}
        \hline
        \textbf{Description}            & \textbf{Length of Compounding Period} & \textbf{Rate per Period} & \textbf{\$1 grows to what after 3 years?} \\
        \hline
        1. 10\% compounded semiannually &                                       &                          &                                           \\
        \hline
        2. 6\% compounded quarterly     &                                       &                          &                                           \\
        \hline
        3. 9\% compounded monthly       &                                       &                          &                                           \\
        \hline
        4. 4\% compounded daily         &                                       &                          &                                           \\
        \hline
    \end{tabular}
\end{table}

\vspace{0.5cm}
\ProblemSection{For questions 5-8, assume 100\% annual interest.  In one year,
    \$1 grows to what if interest is compounded}

\begin{questions}
    \setcounter{question}{4}   % Start question numbering at 5
    \question Annually?
    \vspace{.7cm}
    \question Monthly?
    \vspace{.7cm}
    \question Daily?
    \vspace{.7cm}
    \question Continuously? (Hint - it's $\lim_{x \to \infty} (1+\frac{1}{x})^x$)
\end{questions}

\end{document}
